
\clearpage
\section{实验总结}

本次实验围绕 Cipherfour 算法的差分传播与线性传播分析展开，主要完成了 S 盒差分分布表（DDT）和线性分布表（LAT）的计算、四种传播约束模型的构造与验证，以及六轮差分路线和线性路线的自动化搜索。通过将 S 盒传播关系转化为有限二进制点集，并进一步使用线性约束对其进行刻画，实现了从密码学传播性质到约束求解模型的转换。

在任务一中，分别使用 One-hot 法、逐点排除法、Espresso 逻辑化简法及凸包法构造 S 盒传播模型，并通过 Gurobi 枚举全部二进制可行解。实验结果表明，四种方法均能够准确描述 DDT 的 $72$ 个合法传播点和 LAT 的 $122$ 个合法传播点，但不同方法在辅助变量数量与约束规模方面存在明显差异。其中，One-hot 法的建模形式较为直接，逐点排除法的正确性容易验证，而逻辑化简和凸包约简能够在不改变二进制可行解集合的前提下有效减少约束数量。这说明，在保证模型正确性的基础上，进一步优化约束表示是提高自动化分析效率的重要环节。

在任务二中，利用课程实验平台完成了 Cipherfour 的六轮最优差分路线与线性路线搜索，并结合任务一计算得到的 DDT 和 LAT 对结果进行独立核验。差分分析得到的最优路线概率为 $2^{-6}$，与理论计算完全一致。线性分析则发现，平台报告的整数化权重与根据 LAT 计算得到的精确相关度存在差异。通过进一步检查平台生成的约束文件，分析了这一差异与局部传播权重表示方式之间的关系，也认识到自动化搜索结果的最优性需要在具体模型及目标函数的定义下进行判断。

总体而言，本次实验加深了对差分分析、线性分析及 MILP 建模方法之间联系的理解，也使我更加熟悉了从传播点集构造、线性约束生成到求解器验证的完整流程。特别是在自动化平台的配置与结果分析过程中，进一步认识到准确描述算法结构、统一比特位序以及独立验证求解结果的重要性。
